\documentclass[man,floatsintext,donotrepeattitle]{apa7}

% ---------------------------------------------------------------------------
% Draft manuscript, APA7 format. Target: Journal of Gambling Studies or
% Journal of Behavioral Addictions (both accept APA-style manuscripts on
% submission; final typesetting is done by the journal, so this is a working
% draft, not a camera-ready file). Filled in with the real numbers from
% analysis/greyhound_report.Rmd as of the date below -- re-pull that report
% before submission, since Study 6 is still an open pilot and its N will move.
% ---------------------------------------------------------------------------

\usepackage[american]{babel}
\usepackage{csquotes}
\usepackage[style=apa,sortcites=true,sorting=nyt,backend=biber]{biblatex}
\DeclareLanguageMapping{american}{american-apa}
\addbibresource{references.bib}

\usepackage{graphicx}
\usepackage{booktabs}
\usepackage{siunitx}

\title{TITLE PLACEHOLDER: Near misses in a dynamic gambling task: a robust
  self-report effect that does not translate into betting behavior}
\shorttitle{NEAR MISSES, CONFIDENCE, AND BETTING}

\author{Omar D. P\'erez\textsuperscript{1} \& Luke Clark\textsuperscript{2}}
\affiliation{%
  \textsuperscript{1}Department of Industrial Engineering, Universidad de Chile\\
  \textsuperscript{2}Centre for Gambling Research, Department of Psychology,
  University of British Columbia}

\authornote{% TODO: fill in before submission
  Correspondence concerning this article should be addressed to Omar D.
  P\'erez, Department of Industrial Engineering, Universidad de Chile.
  Email: TODO.
  \par\medskip
  \textbf{Author note (draft placeholder).} Add funding sources, conflict of
  interest statement, preregistration links (Study 2; Study 6 if a
  confirmatory run is preregistered before submission), and a data/code
  availability statement (this repository).}

\abstract{%
  The near-miss effect -- the finding that narrowly missed gambling outcomes
  are experienced differently from clear losses -- is well established for
  self-report and psychophysiological measures, but it is unclear whether it
  changes what people actually do with their money. Across five online
  studies (total \emph{N} = TODO after exclusions; TODO before exclusions)
  using realistic, dynamic gambling stimuli (real greyhound race video,
  rather than the static slot-machine reels used in most prior work),
  participants picked a dog to follow on each race and rated their
  confidence, pleasure, motivation to continue, and perceived luck. The
  near-miss effect on self-report replicated in every study (motivation,
  pleasure, and luck all higher after a near miss than a clear loss;
  confidence unaffected, as it should be if trap assignment was not
  leaking). In two studies that added a real bet-size decision before each
  race (Study 3: 5 discrete Cambridge Gambling Task-style stakes; Study 4: a
  continuous 0--100\% slider), bet size showed no near-miss effect in either
  study, and three statistical approaches with different assumptions about
  the skewed betting data agreed the null was real, not a range or
  measurement artifact. A cross-study comparison suggested that pre-race
  confidence fell as studies added betting; a final, preregistered-in-spirit
  study (Study 6, \emph{N} = 120) tested this directly with a clean
  between-subjects manipulation of bet availability (holding everything else
  about the task constant) and confirmed it: participants who could bet
  reported lower confidence than participants who never could, regardless of
  outcome (\emph{d} = --0.47, one-tailed \emph{p} = .006). The near-miss
  effect on self-report is therefore highly robust, but dissociates from a
  concrete financial decision (bet size); separately, the mere availability
  of a bet -- independent of any near miss -- has a general suppressive
  effect on confidence that has not, to our knowledge, been demonstrated
  before in this direct a comparison.
  \\
  \textbf{Keywords:} near-miss effect, gambling, confidence, illusion of
  control, behavioral persistence, online experiment}

\begin{document}
\maketitle

\section{Introduction}

% ---------------------------------------------------------------------------
% DRAFTING NOTES (delete before submission):
% - This section is a real first draft, not a skeleton, but every citation
%   needs to be checked against the source (see references.bib header).
% - Paragraph order: (1) what the near-miss effect is and why it matters,
%   (2) the standard paradigm and its limits (static, discrete-trial slot
%   machine stimuli), (3) what a dynamic/naturalistic stimulus adds, (4) the
%   open question this project actually answers -- does the effect extend to
%   money-relevant behavior, and does merely being able to bet change how
%   confident people feel, (5) explicit aims/hypotheses per study.
% ---------------------------------------------------------------------------

A near miss in gambling -- a losing outcome that comes close to a win, such
as two matching symbols and a third one line away on a slot machine reel --
is not treated by most gamblers as an ordinary loss. \textcite{reid1986}
first described near-misses as functioning like an intermittent
reinforcement signal despite carrying no informational value about future
outcomes, and \textcite{clark2009} showed that near-misses recruit
win-related brain circuitry (ventral striatum, anterior insula) and increase
self-reported motivation to continue playing, relative to a clear loss, even
though both outcomes pay nothing. This near-miss effect has since been
replicated many times, chiefly using slot-machine or roulette-style
paradigms in which the ``near miss'' is a static, discrete outcome shown for
a fixed duration.

That paradigm has a limitation for generalizing to real gambling: most
naturalistic gambling events unfold over time, and a near miss is often only
apparent from \emph{how} an outcome unfolds, not just the final frame. A
horse or greyhound that leads for most of a race and is caught right at the
line, one that trails throughout and then closes fast without quite getting
there, and one that is never really in contention produce the same binary
outcome (loss) but plausibly very different subjective experiences.
Real races also make available a trajectory dimension -- gaining, fading, or
staying level -- that a static slot-machine reel cannot. TODO: cite any
existing work on dynamic/naturalistic near-miss stimuli here if it exists;
otherwise state explicitly that this is (to our knowledge) the first study
to use real, unedited race footage and to manipulate trajectory this way.

A second open question is whether the near-miss effect, robustly shown on
self-report, extends to a concrete financial decision. Demonstrating that
near-misses make people say they want to keep playing is not the same as
demonstrating that near-misses make people bet more. \textcite{studer2015}
found that recent outcomes shifted confidence and separately shifted
betting, but confidence and betting were measured in different experiments,
with different samples of participants, explicitly cautioning against using
one as a proxy for the other. \textcite{knoxinkster1968} showed that real
racetrack bettors reported higher confidence in their horse
\emph{after} placing a bet than bettors queueing to bet on the same horse
before placing it -- a postdecision effect in the opposite direction from a
simple stakes-lower-confidence account, and one 1991 replication attempt
found no moderation of that effect by bet size. Neither result directly
answers whether \emph{having a bet available at all}, independent of a near
miss and independent of any single decision, changes confidence -- which is
the question the present studies converged on.

We report five studies using real 2013 UK greyhound race footage. Studies 1
and 2 establish and replicate the near-miss effect on self-report in this
naturalistic paradigm, and Study 2 (preregistered) additionally manipulates
the \emph{trajectory} of a near miss (catching up late, leading and then
caught, or staying close throughout) to ask whether the shape of a near miss
matters, not just its presence. Studies 3 and 4 add a real bet-size decision
before each race to test whether the near-miss effect extends to a financial
stake. Study 6 follows up an unplanned observation from Studies 1--4 -- that
average confidence was lower in the two studies that included betting -- with
a clean, randomized, between-subjects test of whether bet availability
itself lowers confidence.

\section{General Method}

% Shared elements across all studies -- write each Study's Method section to
% reference back to this one and only describe what differs.

\subsection{Participants and recruitment}
Participants were recruited on Prolific and completed the task on Pavlovia.
Studies 1--5 participants were excluded from later studies in the series
(Study 6 excluded participants from Studies 1--5). TODO: base pay and
exact eligibility criteria (locale, approval rate) per study.

\subsection{Stimuli}
Race videos were drawn from a fixed pool of real, unedited 2013 UK greyhound
races (12 ``close'' races used for the near-miss/narrow-win manipulation, 4
per trajectory, plus 11 ``clear'' races used for the clear-win/clear-loss
conditions). On every trial, participants chose one of the dogs before the
race and then watched the video. \emph{Trajectory} was coded from the
losing dog's perspective on the home stretch as \emph{catch-up} (gaining
ground but running out of track), \emph{fall-back} (leading, then passed
before the line), or \emph{stable} (close throughout, gap never changing).

\subsection{Conditions}
Each trial ended in one of four conditions from the participant's chosen
dog's perspective: clear win (CW), narrow win (NW), near miss (NM), or clear
loss (CL). Role assignment (which dog is followed, and in the close races,
whether that dog's near/narrow role is a win or a loss) was counterbalanced
by rotation across the sample (a deterministic assignment keyed to
participant order/ID), not drawn at random, so that every clip serves in
both roles an equal number of times across the sample -- with the single
exception of Study 1, which used random assignment before this rotation
scheme was adopted (see Study 1).

\subsection{Measures}
Before each race, participants rated their confidence (0--100 slider) that
their chosen dog would win. After the outcome, they rated how pleased they
felt, how motivated they were to keep playing, and how lucky they felt
(0--100 sliders each). Confidence is taken \emph{before} the outcome and
should not differ across post-hoc outcome categories; it therefore serves
throughout as a manipulation-check that role/trap assignment is not leaking
information. Every study also administered the Problem Gambling Severity
Index \parencite{ferriswynne2001} once, at the end.

\subsection{Data quality and exclusions}
Common exclusion criteria across studies: incomplete sessions, failing an
attention/slider check, straight-lining (zero variance) on any self-report
measure, and (from Study 2 onward) failing an outcome-check question asked
immediately after the race and before the result screen. Exact criteria and
counts are reported per study below.

\subsection{Analysis}
Unless noted otherwise, the near-miss effect on a given measure is tested as
a paired-samples comparison of near miss (NM) versus clear loss (CL) trials,
averaged within participant. We report Bayes factors (JZS prior, $r =
.707$) alongside every $p$-value, so that a null result can be described as
evidence \emph{for} the null ($BF_{01} > 3$) rather than only as a failure
to reject it.

\section{Study 1: The Near-Miss Effect in a Dynamic Task}

\subsection{Method}
Trap assignment (which dog a participant follows, and thus which close races
become near misses versus narrow wins) was drawn at random rather than by
rotation, so per-clip role counts are somewhat uneven (reported as a
manipulation-check table in the full analysis report). Twenty trials per
participant, five of each condition (CW, NW, NM, CL). \emph{N} = 44
analyzed of 51 collected; exclusions: incomplete sessions, failed attention
check, and straight-lining.

\subsection{Results}
The near-miss effect replicated on every self-report measure except
confidence, as expected: motivation was $+9.8$ points higher after a near
miss than a clear loss ($p = 3\times10^{-6}$). Pleasure and luck showed the
same pattern (full test statistics in the online supplement / analysis
report). Confidence, rated before the outcome, was flat between NM and CL,
confirming that trap assignment was not leaking.

\section{Study 2: Does the Shape of the Near Miss Matter?}

\subsection{Method}
Preregistered on AsPredicted. \emph{N} = 102 analyzed of 181 collected.
Role assignment now used the rotation scheme described in General Method,
and an outcome-check question (asked twice, immediately post-race) was
added to catch inattentive responding.

\subsection{Results}
The near-miss effect on motivation replicated ($+10.1$, $p = 2\times10^{-13}$).
Two planned contrasts on the near-miss trials tested whether trajectory
(catch-up, fall-back, stable) modulates the effect. Motivation was flat
across trajectory, but pleasure was not: a catch-up near miss felt
distinctly less unpleasant than either a fall-back or a stable near miss,
which did not differ from each other -- the shape of a near miss changes how
it \emph{feels} without changing how much people want to keep playing.
Confidence again stayed flat across trajectory, as required.

An exploratory analysis asked whether Problem Gambling Severity Index (PGSI)
score moderates the catch-up advantage on pleasure. The advantage was
roughly three times larger in the elevated-PGSI band ($\geq 3$) than the low
band, and the interaction reached significance when PGSI was split at the
standard clinical cutoff, but not when PGSI was entered as a continuous
predictor -- treated as suggestive, not confirmatory, moderation evidence
(about half the sample scored zero on the PGSI, limiting power for this
test, as anticipated in advance).

\section{Study 3: Does the Near-Miss Effect Extend to Bet Size?}

\subsection{Method}
Same races, trajectory coding, and rotation scheme as Study 2. A new
decision was added before every race: how much of the current balance to
stake, chosen from five fixed options modeled on the Cambridge Gambling Task
(5\%, 25\%, 50\%, 75\%, 95\%). Starting balance 100 points, floor 20;
real payment was decoupled from the literal final balance (a fixed
component plus a capped variable bonus) so no participant's actual earnings
depended on the arithmetic of proportional betting. Unregistered pilot.
\emph{N} = 40 analyzed.

\subsection{Results}
Self-report replicated the near-miss effect as in Studies 1--2 ($+6.5$ on
motivation, $p = .008$). Bet size on the race \emph{following} a near miss
versus a clear loss showed no difference. However, 73\% of all bets fell on
the two lowest stake options, leaving little room for any condition
difference to appear -- a likely floor/range problem rather than evidence
against a behavioral near-miss effect, motivating Study 4.

\section{Study 4: A Continuous Bet-Size Measure}

\subsection{Method}
Identical to Study 3 except the five fixed stakes were replaced with a
continuous 0--100\% slider. Unregistered pilot. \emph{N} = 39 analyzed of
75 collected.

\subsection{Results}
The bet-size distribution now used the full range, with a genuine tail out
to 90--100\% in every condition -- the range problem from Study 3 was fixed.
The near-miss effect on bet size was nonetheless still flat, and this null
held up under three tests that make different assumptions about the
right-skewed betting data: a paired $t$-test on the mean ($p = .695$,
$BF_{01} = 5.4$), a Wilcoxon signed-rank test on the median ($p = .385$),
and a paired $t$-test on logit-transformed stakes ($p = .926$, $BF_{01} =
5.8$). All three agree this is a real null, not an artifact of the wrong
test for skewed data.

Self-report replicated again, though more weakly than in Studies 1--3:
motivation was $+3.0$ points higher after a near miss ($p = .1$), the only
one of the five studies in which this specific test did not reach
significance, though the direction was unchanged.

\section{An Unplanned Cross-Study Observation}

Pre-race confidence, pooled across all conditions, differed noticeably
across the four studies run so far: 51.8 (Study 1), 46.7 (Study 2), 40.1
(Study 3), and 35.3 (Study 4) (all on a 0--100 scale). The two studies with
a real bet-size decision (3 and 4) had the two lowest confidence levels.
This comparison is confounded with several other things that changed across
these studies over time (consent screen, payment complexity, an added
decision step, non-overlapping recruitment waves), so it does not on its own
establish that betting availability lowers confidence -- but it generated
that hypothesis, tested cleanly in Study 6.

\section{Study 6: Does Bet Availability Itself Lower Confidence?}

\subsection{Method}
Same race pool, rotation scheme, condition counts, and questions as Study 4.
Participants were randomly assigned once, before consent, to one of two
whole-session conditions: a \emph{betting} group (Study 4's 0--100\% slider,
on every race) or a \emph{no-betting} group (never offered a bet, and never
shown a balance). The direction of the hypothesis (betting lowers
confidence) was fixed before this data existed, from the pattern above, so
the primary test is one-tailed. \emph{N} = 120 analyzed (65 betting, 55
no-betting).

% TODO before submission: decide and state whether this N is the final,
% locked sample (stop-collection decision) or an interim look, and whether
% a separate confirmatory preregistered replication is planned/run. As of
% this draft the target power analysis called for roughly 45-100 per arm
% depending on assumed effect size; check the current analysis report for
% the up-to-date number before writing this paragraph for real.

\subsection{Results}
Participants in the betting group reported lower average pre-race
confidence (\emph{M} = 40.2, \emph{SD} = 20.9) than participants in the
no-betting group (\emph{M} = 49.6, \emph{SD} = 19.0), $t(117.3) = -2.58$,
one-tailed $p = .006$, $d = -0.47$, $BF_{10} = 7.1$ -- moderate evidence for
the effect over the null. The classic near-miss effect on motivation
replicated within this sample as well, pooled across both groups
($+10.3$, $p = 5\times10^{-10}$), confirming the manipulation and task were
behaving normally in both arms.

\section{General Discussion}

Across five independent samples and three variants of this dynamic,
naturalistic gambling task, the near-miss effect on self-reported
motivation, pleasure, and perceived luck was remarkably consistent: near
misses made people feel and say they wanted to continue more than a clear
loss, while confidence -- rated before the outcome was known -- stayed flat,
exactly as a valid manipulation should produce. This replicates the
foundational self-report finding \parencite{clark2009} in a task built from
real, dynamic race footage rather than static slot-machine outcomes, and
extends it with a trajectory dimension not available in prior paradigms:
Study 2 showed that a near miss that involves a late charge (catch-up) feels
less bad than one where the lead was surrendered (fall-back) or where the
gap never changed (stable), even though motivation to continue does not
depend on this distinction.

That self-report robustness makes the second main finding more striking:
across two independently piloted bet-size measures using different scales
(five fixed Cambridge Gambling Task-style stakes in Study 3; a continuous
slider in Study 4), the near-miss effect on bet size was flat, and in Study
4 this null was confirmed under three different statistical approaches that
do not share the same assumptions about the skewed betting data. Wanting to
keep playing, as measured by a rating, did not translate into staking more
money on the next race. This is a genuine dissociation, not a failure to
detect an effect that is there: Study 4's histogram shows real variance and
a genuine tail to the top of the scale in every condition, ruling out the
range/floor explanation that could be raised against Study 3.

The third finding was not part of the original design. An unplanned
observation that average confidence fell as studies added a bet-size
decision motivated Study 6, a clean, randomized, between-subjects test:
merely being able to bet -- with the near-miss manipulation, race pool, and
every other feature of the task held constant -- lowered confidence relative
to a group that could never bet, regardless of what happened in any given
race. To our knowledge this is a more direct test of whether stakes being
available changes confidence than is available in the existing literature:
\textcite{knoxinkster1968} compared confidence before versus after a single
bet decision, not a between-groups comparison of ever-available versus
never-available stakes, and found postdecision confidence \emph{increased}
after commitment -- the opposite direction from the present result, though
answering a different question. \textcite{studer2015} measured confidence
and betting in separate experiments with different samples, explicitly
cautioning against treating one as a stand-in for the other. Whether the
mechanism here is a general conservatism induced by real stakes, an
attentional or cognitive-load cost of the extra decision, or something else
is not resolved by this design and is worth its own follow-up.

\subsection{Limitations}
Studies 3, 4, and 6 were unregistered pilots (Study 6's hypothesis direction
was fixed in advance from the cross-study pattern, which licenses a
one-tailed test, but the study as a whole was not formally preregistered).
Study 6's current sample, while sufficient to detect the observed effect
size, is smaller than would be wanted for a fully powered confirmatory claim
if the true effect is closer to the small-to-medium end of plausible values;
TODO -- update this paragraph with the final, locked sample size and state
plainly whether a larger confirmatory run was completed before submission.
All self-report measures are single-item 0--100 sliders rather than
validated multi-item scales. The Prolific sample, like most online panels,
is not enriched for problem gambling (roughly half of participants scored
zero on the PGSI in every study), limiting the power of the PGSI moderation
analyses, consistent with prior cautions about testing this kind of
moderation in general-population online samples.

\subsection{Future directions}
Because bet size did not carry the near-miss effect, a natural next
behavioral measure is \emph{persistence} -- whether someone keeps playing at
all -- rather than how much they stake when they do. Adapting this task's
existing (currently unused) cash-out mechanic into a repeated, per-trial
stop-or-continue offer would give a persistence measure tied to the specific
condition that just occurred, within the same repeated-measures structure
used throughout this project. A sample enriched for problem gambling would
give the PGSI-moderation questions raised in Study 2 real power. Finally, a
dedicated cashout-manipulation study could ask whether the near-miss effect
depends on the availability of a cashout decision, whether trajectory
moderates that dependency, and whether a near miss's trajectory itself
changes cashout behavior -- combining the open questions from Studies 2 and
6 with the persistence measure above.

\printbibliography

\end{document}
